\documentclass[10pt,a4paper]{amsart}
\usepackage[margin=19mm]{geometry}
\usepackage{amsmath,amssymb,mathtools}
\usepackage[scr=rsfs,cal=euler]{mathalfa}
\usepackage{xfrac}
\usepackage[unicode,hidelinks,bookmarks=false]{hyperref}
\newcommand{\ie}{i.e.\ }
\newcommand{\cf}{cf.\ }
\newcommand{\resp}{respectively\ }
\newcommand{\Subsection}{\subsection}
\providecommand{\nobreakdash}{\nobreak\hbox{-}}
\setlength{\emergencystretch}{2em}
\multlinegap0pt
\usepackage{bm}
\usepackage{bbm}
\usepackage{dsfont}
\usepackage{xspace}
\usepackage{cancel}
\usepackage{mathtools}
\usepackage{amsthm}
\makeatletter
\@ifundefined{lem}{\newtheorem{lem}{Lemma}}{}
\@ifundefined{thm}{\newtheorem{thm}{Theorem}}{}
\makeatother
\title[Inverse elastic obstacle scattering problems by monotonicity]{Inverse elastic obstacle scattering problems by monotonicity method}
\author{Mengjiao Bai, Huaian Diao, Weisheng Zhou}
\date{Source: arXiv:2506.04655v1 (CC BY 4.0)}
\begin{document}
\maketitle

\noindent\emph{Source and licence.} Inverse elastic obstacle scattering problems by monotonicity method. Version arXiv:2506.04655v1 (CC BY 4.0), distributed under the Creative Commons Attribution 4.0 International licence, as recorded in the arXiv licence metadata for this exact version and shown on the abstract page https://arxiv.org/abs/2506.04655v1. Full rights notice: licenses/arxiv-2506.04655-CC-BY-4.0.txt.\par\smallskip

\noindent\textit{Editorial scope.} Counted unit: the Thm whose source label is \texttt{LocialWaveFounc} in the article (source0.tex lines 480--485, source label \texttt{LocialWaveFounc}), delivered together with the article's own complete original proof of it (source0.tex lines 487--564, one \texttt{proof} environment of 78 lines). No mathematics is written, repaired or re-derived: statement and proof are the article's own text line for line. Required context retained uncounted: eq:RestriH (lines 416--418); RangeHG (lines 420--425); Inquallity (lines 464--470); Dim (lines 472--478). Every definition, display and label of those lines is retained unchanged. Omitted from this package and not redistributed: 1--415, 419--419, 426--463, 471--471, 479--479, 486--486, 565--856. Nothing inside a retained excerpt is omitted or altered: every retained body is delivered exactly as the article's lines are written, so no omission receipt of any kind accompanies this package. Citations: the retained excerpts carry no citation command at all, so no bibliography entry of the article is transcribed and no reference is redistributed. Reference adaptation: none was needed and none was made; every reference inside the delivered unit resolves inside it, so no unresolved reference or citation is produced. Printed slips kept literal: no printed word, symbol, hypothesis or number of the counted statement or of its proof was altered. Layout and class: the article's own class is \texttt{amsart} with packages including graphicx, amssymb, epstopdf, geometry, hyperref, booktabs, array, paralist, verbatim, subfig, tabularx, amsmath, amsfonts, amsthm; the wrapper is the lane's pinned \texttt{amsart} editorial wrapper at 10pt on A4, which restates the theorem-like environments the retained text needs and adds the article's own macro definitions that the retained text needs (none), with extra wrapper packages (bm, bbm, dsfont, xspace, cancel, mathtools). The delivered numbering is the unit's own. Assets: no retained span contains a figure environment, a graphics-inclusion command, a PDF-page inclusion, a table or a TikZ picture; no image file is copied into this package, the package declares no asset dependency and no image file is redistributed with it. Licence: version arXiv:2506.04655v1 of this article is distributed under the Creative Commons Attribution 4.0 International licence, as recorded in the arXiv licence metadata for this exact version and shown on the abstract page https://arxiv.org/abs/2506.04655v1. A full rights notice is delivered at licenses/arxiv-2506.04655-CC-BY-4.0.txt. Characters: every retained excerpt is pure ASCII, so the delivered engine, XeLaTeX with its default Latin Modern text font, reports no missing character.
	\begin{align}\label{eq:RestriH}
		\mathbf H^*_{\tau}=\mathbf H^*_B \mathbf R^*_{\tau}=\sqrt{8\pi\omega}\mathbf{G_B S_B}  \mathbf R^*_{\tau}.
	\end{align}

	\begin{thm}\label{RangeHG}
		Let $B, D \subset \mathbb R^2$ are open and Lipschitz bounded domain, if $B \nsubseteq D$ ,let $\tau \subseteq \partial B \setminus \overline{D} $ is relatively open and $\mathbb R^2 \setminus \left( \overline{\tau \cup D }\right) $ is connected, then we have
			\begin{align}
				R\left( \mathbf H^*_{\tau} \right) \cap R\left( \mathbf G\right) =\left\lbrace 0\right\rbrace .
			\end{align}
	\end{thm}

\begin{lem}\label{Inquallity}%\cite{AG20} 
	Assume $X, Y, Z$ be Hilbert spaces, let $A_1:X \to Y$ and $A_2: X \to Z$ are linear operators, then the following two statements are equivalent
	\begin{itemize}
		\item [(1).] There exists a constant $C>0$ such that $||A_1 x||_Y \leq C ||A_2 x||_Z, \quad \forall x \in X$.
		\item [(2).] $R \left( A^*_1 \right) \subseteq R \left( A^*_2 \right) $.
	\end{itemize}
\end{lem}

\begin{lem}\label{Dim}%\cite{AG20} 
	Assume $V, Z_1, Z_2 $ be subspaces of a vector space $Z$, if 
	$$
	Z_1 \cap Z_2= \left\lbrace 0 \right\rbrace \quad and \quad Z_1 \subseteq Z_2 +V,
	$$
	then $\rm {dim}\left(  Z_1 \right) \leqslant \rm {dim}\left(  V \right) $.
\end{lem}

\begin{thm}\label{LocialWaveFounc}
 Let $B, D  \subseteq \mathbb R^2$ is open and Lipschitz bounded such that $\mathbb R^2 \backslash \overline {D}$ is connected. Assume $B \nsubseteq D$, then for any finite-dimensional subspace $V \subseteq \left[L^{2}\left(\mathbb S\right) \right] ^2 $ , there exists a sequence $\left\lbrace  \mathbf h_n \right\rbrace  \subseteq  V^{\perp} $ such that
 $$
 \left\|\mathbf H_B \mathbf h_n\right\|_{\left[H^{1/2}\left(\partial B\right) \right] ^2 } \to \infty \quad and \quad \left\| \mathbf G^* \mathbf h_n\right\|_{\left[H^{-1/2}\left(\partial D\right) \right] ^2 } \to 0, \qquad n\to \infty.
 $$
\end{thm}

\begin{proof}
Let $B, D  \subseteq \mathbb R^2$ is open and Lipschitz bounded such that $\mathbb R^2 \backslash \overline {D}$ is connected. Assume $B \nsubseteq D$, let $V \subseteq \left[L^{2}\left(\mathbb S\right) \right] ^2 $ is a finite-dimensional subspace, therefore, the orthogonal projection onto V is well-defined, and we denote it as $\mathbf P :\left[L^{2}\left(\mathbb S\right) \right] ^2 \to \left[L^{2}\left(\mathbb S\right) \right] ^2$. 

Since $B \nsubseteq D $ and $\mathbb R^2 \backslash \overline {D}$ is connected, there exists a relatively open $\tau \subseteq \partial B\backslash \overline{D}$ such that $\mathbb R^2 \backslash \overline {\left( \tau \cup D\right) }$ is connected, from Theorem \ref{RangeHG} we have
\begin{align}\label{eq:RangeHtG}
R\left( \mathbf H^*_{\tau} \right) \cap R\left( \mathbf G\right) =\left\lbrace 0\right\rbrace .
\end{align}

Now, we turn our attention to equation \eqref{eq:RestriH}. Without loss of generality, we assume that $\omega ^2$ is not an eigenvalue of $ \triangle ^*$ in $B$, therefore, both $\mathbf S _B$ (single-layer potential operator on $\partial B$ ) and $\mathbf G _B$ are injective. Moreover, the range of the extension operator $\mathbf R^*_{\tau} $ is infinite-dimensional, which implies that $ R \left( \mathbf H^*_{\tau} \right) = R \left( \sqrt{8\pi\omega}\mathbf{G_B S_B}  \mathbf R^*_{\tau}   \right) $ is infinite-dimensional. Therefore, according to lemma \ref{Dim} and \eqref{eq:RangeHtG} , we can obtain
$$
R \left( \mathbf H^*_{\tau} \right) \nsubseteq R\left( \mathbf G\right) +V=  R \left( \left[ \mathbf G \quad \mathbf P \right] \right).
$$
Hence, utilizing the Lemma \ref{Inquallity}, we can deduce that there does not exist a constant $C > 0$ such that
\begin{equation}\begin{aligned}\label{eq:inver}
|| \mathbf H_{\tau} \mathbf g ||^2_{\left[H^{1/2}\left(\tau \right) \right] ^2 }
&\leq C^2 \left\| 
\begin{bmatrix}
	\mathbf G^* \\ \mathbf P
\end{bmatrix}
\mathbf g \right\| ^2_{\left[H^{-1/2}\left(\partial D \right) \right] ^2  \times \left[L^{2}\left(\mathbb S \right) \right] ^2}\\
&=  C^2 \left( \left\| \mathbf G^* \mathbf g \right\|^2_{\left[H^{-1/2}\left(\partial D \right) \right] ^2}+\left\| \mathbf P \mathbf g \right\|^2_{\left[L^{2}\left(\mathbb S \right) \right] ^2} \right) .
\end{aligned}\end{equation}
Since $\mathbf P$ is an orthogonal projection operator, it follows that $\mathbf P$ is a self-adjoint operator, that is, $\mathbf P = \mathbf P ^*$. Therefore, from \eqref{eq:inver} for any $n \in \mathbb N$, there exists a $\mathbf g_n \in \left[L^{2}\left(\mathbb S \right) \right] ^2  $ such that
$$
\left\| \mathbf H_{\tau} \mathbf g_n \right\|^2_{\left[H^{1/2}\left(\tau \right) \right] ^2 } > n^2 \left( \left\| \mathbf G^* \mathbf g_n \right\|^2_{\left[H^{-1/2}\left(\partial D \right) \right] ^2}+\left\| \mathbf P \mathbf g_n \right\|^2_{\left[L^{2}\left(\mathbb S \right) \right] ^2} \right).
$$
For any $n \in \mathbb N$, we denote $M_n:=\left\| \mathbf G^* \mathbf g_n \right\|^2_{\left[H^{-1/2}\left(\partial D \right) \right] ^2}+\left\| \mathbf P \mathbf g_n \right\|^2_{\left[L^{2}\left(\mathbb S \right) \right] ^2}$,  $\tilde{\mathbf g_n}:=\mathbf g_n /\left(\sqrt{n M_n} \right) $. Then on one hand, we have
\begin{equation}\begin{aligned}
\left\| \mathbf H_{\tau} \tilde{\mathbf g_n} \right\|^2_{\left[H^{1/2}\left(\tau \right) \right] ^2 }
&= \frac{1}{n M_n} \left\| \mathbf H_{\tau} \mathbf g_n \right\|^2_{\left[H^{1/2}\left(\tau \right) \right] ^2 } \\
&>\frac{n}{M_n} \left( \left\|  \mathbf G^* \mathbf g_n \right\|^2_{\left[H^{-1/2}\left(\partial D \right) \right] ^2}+\left\|  \mathbf P \mathbf g_n \right\|^2_{\left[L^{2}\left(\mathbb S \right) \right] ^2} \right) \\
&= n.
\end{aligned}\end{equation}
i.e. 
$$
\left\| \mathbf H_{\tau} \tilde{\mathbf g_n} \right\|^2_{\left[H^{1/2}\left(\tau \right) \right] ^2 } > n \to \infty,\quad   n \to \infty .
$$
On the other hand, 
\begin{equation}\begin{aligned}
		\left\| \mathbf G^* \tilde{\mathbf g_n} \right\|^2_{\left[H^{-1/2}\left(\partial D \right) \right] ^2}+\left\| \mathbf P \tilde{\mathbf g_n} \right\|^2_{\left[L^{2}\left(\mathbb S \right) \right] ^2}
		&= \frac{1}{n M_n} \left(\left\| \mathbf G^* \mathbf g_n \right\|^2_{\left[H^{-1/2}\left(\partial D \right) \right] ^2}+\left\| \mathbf P \mathbf g_n \right\|^2_{\left[L^{2}\left(\mathbb S \right) \right] ^2} \right) \\
		&= \frac{1}{n}.
\end{aligned}\end{equation}
Thus, we can obtain
$$
\left\| \mathbf G^* \tilde{\mathbf g_n} \right\|^2_{\left[H^{-1/2}\left(\partial D \right) \right] ^2}+\left\| \mathbf P \tilde{\mathbf g_n} \right\|^2_{\left[L^{2}\left(\mathbb S \right) \right] ^2}= \frac{1}{n} \to 0, \quad n \to \infty .
$$
In summary, when $n \to \infty $, we can obtain
\begin{equation}\begin{aligned}\label{eq:HGPasymptotic}
  \left\| \mathbf H_{\tau} \tilde{\mathbf g_n} \right\|^2_{\left[H^{1/2}\left(\tau \right) \right] ^2 } \to \infty ,\quad \left\| \mathbf G^* \tilde{\mathbf g_n} \right\|^2_{\left[H^{-1/2}\left(\partial D \right) \right] ^2} \to 0,\quad \left\| \mathbf P \tilde{\mathbf g_n} \right\|^2_{\left[L^{2}\left(\mathbb S \right) \right] ^2} \to 0.
 \end{aligned}\end{equation}

For any $n\in \mathbb N$, we define $\mathbf h_n:= \tilde{\mathbf g_n} - \mathbf P \tilde{\mathbf g_n}$, using the triangle inequality, we can obtain
\begin{equation}\begin{aligned}\label{eq:HInequality}
	    \left\| \mathbf H_{\tau} \mathbf h_n \right\|_{\left[H^{1/2}\left(\tau \right) \right]^2 }
		&\geq \left| \left\| \mathbf H_{\tau} \tilde{\mathbf g_n} \right\|_{\left[H^{1/2}\left(\tau \right) \right]^2 } - \left\| \mathbf H_{\tau} \mathbf P \tilde{\mathbf g_n} \right\|_{\left[H^{1/2}\left(\tau \right) \right]^2 } \right |\\
		&\geq \left\| \mathbf H_{\tau} \tilde{\mathbf g_n} \right\|_{\left[H^{1/2}\left(\tau \right) \right]^2 } - \left\| \mathbf H_{\tau} \mathbf P \tilde{\mathbf g_n} \right\|_{\left[H^{1/2}\left(\tau \right) \right]^2 }\\
		&\geq \left\| \mathbf H_{\tau} \tilde{\mathbf g_n} \right\|_{\left[H^{1/2}\left(\tau \right) \right]^2 } - \left\| \mathbf H_{\tau} \right\|_{ \left[L^{2}\left(\mathbb S \right) \right]^2 \to \left[H^{1/2} \left(\tau \right) \right] ^2 } \left\|\mathbf P \tilde{\mathbf g_n} \right\|_{\left[L^{2} \left(\mathbb S \right) \right] ^2}.
\end{aligned}\end{equation}
and
\begin{equation}\begin{aligned}\label{eq:GInequality}
		\left\| \mathbf G^* \mathbf h_n \right\|_{\left[H^{-1/2}\left(\partial D \right) \right]^2 }
		&\leq \left\| \mathbf G^* \tilde{\mathbf g_n }\right\|_{\left[H^{-1/2}\left(\partial D \right) \right]^2 }+ \left\| \mathbf G^* \mathbf P \tilde{\mathbf g_n }\right\|_{\left[H^{-1/2}\left(\partial D \right) \right]^2 } \\
		&\leq \left\| \mathbf G^* \tilde{\mathbf g_n }\right\|_{\left[H^{-1/2}\left(\partial D \right) \right]^2 }+ \left\| \mathbf G^*\right\|_{\left[L^{2} \left(\mathbb S \right) \right] ^2 \to \left[H^{-1/2}\left(\partial D \right) \right]^2} \left\|\mathbf P \tilde{\mathbf g_n }\right\|_{\left[L^{2} \left(\mathbb S \right) \right] ^2}.
\end{aligned}\end{equation}
From \eqref{eq:HGPasymptotic} , \eqref{eq:HInequality} and  \eqref{eq:GInequality} we have $ \left\| \mathbf H_{\tau} \mathbf h_n \right\|_{\left[H^{1/2}\left(\tau \right) \right]^2 } \to \infty $, $\left\| \mathbf G^* \mathbf h_n \right\|_{\left[H^{-1/2}\left(\partial D \right) \right]^2 } \to 0$, when $n \to \infty$. From the definition of the restriction operator $\mathbf R_{\tau}$, we have that $\mathbf H_{\tau}= \mathbf R_{\tau} \mathbf H_{B} $ holds, so we can obtain
\begin{equation}\begin{aligned}
\left\| \mathbf H_{\tau} \mathbf h_n \right\|_{\left[H^{1/2}\left(\tau \right) \right]^2 } 
&=\left\|\mathbf R_{\tau} \mathbf H_{B} \mathbf h_n \right\|_{\left[H^{1/2}\left(\partial B \right) \right]^2 } \\
& \leq \left\|\mathbf R_{\tau}\right\|_{\left[H^{1/2}\left(\partial B \right) \right]^2 \to \left[H^{1/2}\left(\tau \right) \right]^2 } \left\| \mathbf H_{B} \mathbf h_n \right\|_{\left[H^{1/2}\left(\partial B \right) \right]^2 } ,
\end{aligned}\end{equation}
i.e. we have $ \left\| \mathbf H_{B} \mathbf h_n \right\|_{\left[H^{1/2}\left(\partial B \right) \right]^2 } \geq \left\| \mathbf H_{\tau} \mathbf h_n \right\|_{\left[H^{1/2}\left(\tau \right) \right]^2 }  / \left\|\mathbf R_{\tau}\right\|_{\left[H^{1/2}\left(\partial B \right) \right]^2 \to \left[H^{1/2}\left(\tau \right) \right]^2 }$, since when $n \to \infty $,  it follows that $\left\| \mathbf H_{\tau} \mathbf h_n \right\|_{\left[H^{1/2}\left(\tau \right) \right]^2 } \to \infty $ and $\left\|\mathbf R_{\tau}\right\|_{\left[H^{1/2}\left(\partial B \right) \right]^2 \to \left[H^{1/2}\left(\tau \right) \right]^2 }$ is a constant, we have
$$
\left\| \mathbf H_{B} \mathbf h_n \right\|_{\left[H^{1/2}\left(\partial B \right) \right]^2 } \to \infty ,\quad n\to\infty.
$$

Therefore, with the above expression and \eqref{eq:GInequality} , the proof is completed.
\end{proof}

\end{document}
